\chapter{Complexity — Basic}
%%% generated by the blueprint pipeline from TCSlib.Complexity.ClassP.BinaryArith.Basic
\section{Overview}

Declarations of the Lean module \texttt{TCSlib.Complexity.ClassP.BinaryArith.Basic}.

\section{Declarations}

\begin{theorem}[Binary value from the first bit]
\lean{Complexity.BinArith.bitsVal_eq_headD}
\label{Complexity.BinArith.bitsVal_eq_headD}
\leanok
\uses{Complexity.LogProg.bitsVal}
\difficulty{4}
For a finite Boolean word $w$, write $V(w)$ for its binary value, with the first bit
least significant. Then $V(w)=\beta+2V(w')$, where $\beta=0$ if $w$ is empty and otherwise
is $0$ or $1$ according as its first bit is false or true, and $w'$ is obtained by
deleting the first bit (or is empty if $w$ is empty).
\end{theorem}

\begin{theorem}[Value of an appended bit]
\lean{Complexity.BinArith.bitsVal_append_singleton}
\label{Complexity.BinArith.bitsVal_append_singleton}
\leanok
\uses{Complexity.LogProg.bitsVal}
\difficulty{4}
For a finite Boolean list $w$ and a bit $b$, the value of the list obtained by appending
$b$ is the value of $w$ plus $b$ times $2^{|w|}$, where the list value interprets its
first bit as the least significant bit.
\end{theorem}

\begin{theorem}[Zero binary value criterion]
\lean{Complexity.BinArith.bitsVal_eq_zero_iff}
\label{Complexity.BinArith.bitsVal_eq_zero_iff}
\leanok
\uses{Complexity.LogProg.bitsVal}
\difficulty{5}
For every finite Boolean word $w$, its binary value is zero if and only if $w$ contains
no true bit.
\end{theorem}

\begin{theorem}[Binary expansion length bound]
\lean{Complexity.BinArith.length_bits_bitsVal_le}
\label{Complexity.BinArith.length_bits_bitsVal_le}
\leanok
\uses{Complexity.LogProg.bitsVal, Complexity.LogProg.bitsVal_lt_two_pow, Complexity.LogProg.length_bits_le}
\difficulty{1}
For every finite bit string $w$, the length of the binary expansion of its value, with
the first bit interpreted as least significant, is at most the length of $w$.
\end{theorem}

\begin{theorem}[Binary expansion length from a bounded representation]
\lean{Complexity.BinArith.length_bits_le_of_le}
\label{Complexity.BinArith.length_bits_le_of_le}
\leanok
\uses{Complexity.BinArith.length_bits_bitsVal_le, Complexity.LogProg.bitsVal, Complexity.LogProg.length_bits_mono}
\difficulty{3}
Let $n,m$ be natural numbers and let $w$ be a finite bit string whose value is at least
$n$ and whose length is at most $m$. Then the length of the binary expansion of $n$ is
at most $m$.
\end{theorem}

\begin{theorem}[Repeated deletion of the first bit]
\lean{Complexity.BinArith.iterate_drop_one}
\label{Complexity.BinArith.iterate_drop_one}
\leanok
\difficulty{4}
For every bit string $w$ and natural number $i$, deleting the first bit repeatedly $i$
times gives the suffix of $w$ obtained by deleting its first $i$ bits.
\end{theorem}

\begin{theorem}[Polynomial-time nonemptiness test]
\lean{Complexity.BinArith.polyTimeComputable_any}
\label{Complexity.BinArith.polyTimeComputable_any}
\leanok
\uses{Complexity.PolyTimeComputable, Complexity.PolyTimeComputable.comp, Complexity.polyTimeComputable_isNil, Complexity.polyTimeComputable_not, Complexity.polyTimeComputable_transduce, Complexity.transduce}
\difficulty{5}
The function on bit strings that returns the one-bit word $[1]$ if the input contains at
least one $1$, and the one-bit word $[0]$ otherwise, is polynomial-time computable.
\end{theorem}

\begin{definition}[Normalization of a Boolean chunk]
\lean{Complexity.BinArith.normChunk}
\label{Complexity.BinArith.normChunk}
\leanok
Given a list of Boolean values, return its first element as a singleton list if the list
contains any `true` value; otherwise return the empty list.
\end{definition}

\begin{theorem}[Binary normalization by chunks]
\lean{Complexity.BinArith.flatMap_normChunk}
\label{Complexity.BinArith.flatMap_normChunk}
\leanok
\uses{Complexity.BinArith.bitsVal_eq_zero_iff, Complexity.BinArith.normChunk, Complexity.LogProg.bitsVal, Complexity.LogProg.bits_two_mul_add}
\difficulty{6}
Let $w$ be a finite bit string of length at most $N$. For each $i=0,\ldots,N-1$, take
the first bit of the suffix of $w$ beginning at position $i$, if that suffix contains a
$1$, and otherwise take the empty string. Concatenating these strings gives the binary
expansion of the natural number represented by $w$, with the least significant bit
first.
\end{theorem}

\begin{theorem}[Polynomial-Time Binary Normalization]
\lean{Complexity.BinArith.polyTimeComputable_normBits}
\label{Complexity.BinArith.polyTimeComputable_normBits}
\leanok
\uses{Complexity.BinArith.flatMap_normChunk, Complexity.BinArith.iterate_drop_one, Complexity.BinArith.normChunk, Complexity.BinArith.polyTimeComputable_any, Complexity.LogProg.bitsVal, Complexity.PolyTimeComputable, Complexity.polyTimeComputable_const, Complexity.polyTimeComputable_emitIter, Complexity.polyTimeComputable_ite, Complexity.polyTimeComputable_tail, Complexity.polyTimeComputable_take1}
\difficulty{4}
For a bit string $w$, interpret its bits as a binary numeral with the first bit least
significant, and return the binary expansion of that natural number. This normalization
function is polynomial-time computable.
\end{theorem}

\begin{definition}[Final element of a Boolean chunk]
\lean{Complexity.BinArith.lastChunk}
\label{Complexity.BinArith.lastChunk}
\leanok
Given a list of Boolean values, return its first element as a singleton list if the list
has at most one element; otherwise return the empty list.
\end{definition}

\begin{theorem}[Concatenation of last chunks]
\lean{Complexity.BinArith.flatMap_lastChunk}
\label{Complexity.BinArith.flatMap_lastChunk}
\leanok
\uses{Complexity.BinArith.lastChunk}
\difficulty{4}
Let $w$ be a bit string and $N$ a natural number with $|w|\le N$. For each
$i=0,\ldots,N-1$, take the last chunk of the suffix of $w$ beginning at position $i$,
and concatenate these chunks in order. The result is the one-term list containing the
last bit of $w$, or the empty list if $w$ is empty.
\end{theorem}

\begin{theorem}[Polynomial-Time Last-Bit Function]
\lean{Complexity.BinArith.polyTimeComputable_lastBit}
\label{Complexity.BinArith.polyTimeComputable_lastBit}
\leanok
\uses{Complexity.BinArith.flatMap_lastChunk, Complexity.BinArith.iterate_drop_one, Complexity.BinArith.lastChunk, Complexity.PolyTimeComputable, Complexity.PolyTimeComputable.comp, Complexity.polyTimeComputable_const, Complexity.polyTimeComputable_emitIter, Complexity.polyTimeComputable_headD, Complexity.polyTimeComputable_isNil, Complexity.polyTimeComputable_ite, Complexity.polyTimeComputable_tail, Complexity.polyTimeComputable_take1}
\difficulty{5}
The function sending a finite Boolean word to the one-bit word containing its last bit,
and sending the empty word to $[\mathsf{false}]$, is polynomial-time computable.
\end{theorem}

\begin{theorem}[Second-slot projection is polynomial-time computable]
\lean{Complexity.BinArith.polyTimeComputable_slot2}
\label{Complexity.BinArith.polyTimeComputable_slot2}
\leanok
\uses{Complexity.PolyTimeComputable, Complexity.PolyTimeComputable.comp, Complexity.pairFstD, Complexity.pairSndD, Complexity.polyTimeComputable_pairFstD, Complexity.polyTimeComputable_pairSndD}
\difficulty{1}
For each bit string $S$, decode $S$ into its two components, then decode the second
component once more and return the first component of that decoding, using the empty
word whenever a decoding fails. This function is polynomial-time computable.
\end{theorem}

\begin{theorem}[Third-slot projection is polynomial-time computable]
\lean{Complexity.BinArith.polyTimeComputable_slot3}
\label{Complexity.BinArith.polyTimeComputable_slot3}
\leanok
\uses{Complexity.PolyTimeComputable, Complexity.PolyTimeComputable.comp, Complexity.pairSndD, Complexity.polyTimeComputable_pairSndD}
\difficulty{1}
For each bit string $S$, decode $S$ into its two components, then decode the second
component and take its second component, and decode that result once more and take its
second component; use the empty word whenever a decoding fails. This function is
polynomial-time computable.
\end{theorem}

\begin{theorem}[Polynomial-time construction of a prefixed output]
\lean{Complexity.BinArith.polyTimeComputable_cons}
\label{Complexity.BinArith.polyTimeComputable_cons}
\leanok
\uses{Complexity.PolyTimeComputable, Complexity.PolyTimeComputable.append}
\difficulty{1}
Let $p$ be a Boolean-valued function on bit strings and $f$ a bit-string-valued
function. If the functions $x\mapsto[p(x)]$ and $f$ are polynomial-time computable, then
the function $x\mapsto p(x)::f(x)$, which prefixes the output of $f$ with the bit
$p(x)$, is polynomial-time computable.
\end{theorem}

\begin{theorem}[Polynomial-time iterated projection]
\lean{Complexity.BinArith.polyTimeComputable_slot4}
\label{Complexity.BinArith.polyTimeComputable_slot4}
\leanok
\uses{Complexity.BinArith.polyTimeComputable_slot3, Complexity.PolyTimeComputable, Complexity.PolyTimeComputable.comp, Complexity.pairSndD, Complexity.polyTimeComputable_pairSndD}
\difficulty{1}
The function on bit strings that applies the total second projection three times is
polynomial-time computable.
\end{theorem}

\begin{theorem}[Polynomial-time bitwise exclusive OR]
\lean{Complexity.BinArith.polyTimeComputable_xor}
\label{Complexity.BinArith.polyTimeComputable_xor}
\leanok
\uses{Complexity.PolyTimeComputable, Complexity.polyTimeComputable_ite, Complexity.polyTimeComputable_not}
\difficulty{4}
Let $p$ and $q$ map bit strings to bits. If the functions $x\mapsto[p(x)]$ and
$x\mapsto[q(x)]$ are polynomial-time computable, then the function
$x\mapsto[p(x)\mathbin{\mathrm{xor}}q(x)]$ is polynomial-time computable.
\end{theorem}

\begin{theorem}[Polynomial-time computability of dropping a prefix]
\lean{Complexity.BinArith.polyTimeComputable_drop}
\label{Complexity.BinArith.polyTimeComputable_drop}
\leanok
\uses{Complexity.PolyTimeComputable, Complexity.PolyTimeComputable.comp, Complexity.polyTimeComputable_id, Complexity.polyTimeComputable_tail}
\difficulty{4}
For every natural number $i$, the function on bit strings that deletes the first $i$
bits is polynomial-time computable.
\end{theorem}

\begin{theorem}[Value after dropping the first bit]
\lean{Complexity.BinArith.bitsVal_drop_one}
\label{Complexity.BinArith.bitsVal_drop_one}
\leanok
\uses{Complexity.BinArith.bitsVal_eq_headD, Complexity.LogProg.bitsVal}
\difficulty{3}
For every finite Boolean list $w$, the value of the list obtained by dropping its first
bit is the integer quotient of the value of $w$ by $2$.
\end{theorem}

\begin{theorem}[Binary value modulo two]
\lean{Complexity.BinArith.bitsVal_mod_two}
\label{Complexity.BinArith.bitsVal_mod_two}
\leanok
\uses{Complexity.BinArith.bitsVal_eq_headD, Complexity.LogProg.bitsVal}
\difficulty{3}
For every finite Boolean word $w$, the remainder of its binary value upon division by
$2$ is $\operatorname{val}(w_0)$, where $w_0$ is the first bit, taken to be false when
$w$ is empty.
\end{theorem}

\begin{definition}[Finite padded zip of Boolean lists]
\lean{Complexity.BinArith.padZip}
\label{Complexity.BinArith.padZip}
\leanok
Given two lists of Boolean values and a natural number $N$, return a list of $N$ ordered
pairs, pairing the current first elements at each step and using `false` whenever either
list is exhausted. The lists are advanced by one position at each step.
\end{definition}

\begin{theorem}[Length of a padded zip]
\lean{Complexity.BinArith.length_padZip}
\label{Complexity.BinArith.length_padZip}
\leanok
\uses{Complexity.BinArith.padZip}
\difficulty{4}
For any two bit strings $a,b$ and any natural number $N$, the list formed by taking $N$
aligned pairs of successive bits from $a$ and $b$, padding a missing bit with $0$, has
length $N$.
\end{theorem}

\begin{definition}[Serial output from a state machine]
\lean{Complexity.BinArith.serialOut}
\label{Complexity.BinArith.serialOut}
\leanok
Let $sigma$ be a state space, let $delta$ update the state from a state and two
Boolean inputs, and let $o$ produce a Boolean output from a state and two Boolean
inputs. Given an initial state and a finite list of input pairs, return the list of
outputs obtained by applying $o$ at each pair using the current state, then updating the
state with $delta$ before processing the next pair.
\end{definition}

\begin{definition}[Final state after serial inputs]
\lean{Complexity.BinArith.serialState}
\label{Complexity.BinArith.serialState}
\leanok
Let $sigma$ be a state space and let $delta$ update the state from a state and two
Boolean inputs. Given an initial state and a finite list of input pairs, return the
state obtained by applying $delta$ successively to each pair, starting from the initial
state.
\end{definition}

\begin{theorem}[Length of serial output]
\lean{Complexity.BinArith.length_serialOut}
\label{Complexity.BinArith.length_serialOut}
\leanok
\uses{Complexity.BinArith.serialOut}
\difficulty{4}
For any set of states, transition function, output function, starting state, and finite
list of bit pairs, the serial output produced by processing the list has the same length
as the list.
\end{theorem}

\begin{theorem}[Final Output Equals Final State]
\lean{Complexity.BinArith.getLastD_serialOut_self}
\label{Complexity.BinArith.getLastD_serialOut_self}
\leanok
\uses{Complexity.BinArith.serialOut, Complexity.BinArith.serialState}
\difficulty{5}
Let $\delta$ be a Boolean state transition function that takes a state and two input bits,
and let $s$ be an initial Boolean state. For any nonempty finite sequence of pairs of
Boolean inputs, if each output bit is computed by applying $\delta$ to the current state and
the next input pair, then the last output bit equals the state obtained after processing
the entire sequence.
\end{theorem}

\begin{definition}[Paired Boolean update]
\lean{Complexity.BinArith.pairδ}
\leanok
Let $\sigma$ be a type and let $\delta : \sigma \times \{0,1\} \times \{0,1\} \to
\sigma$. Define a state update that, when no bit is pending, stores the input bit and
leaves the state unchanged; when a bit $x$ is pending and the next input is $y$, clears
the pending bit and updates the state to $\delta(s,x,y)$.
\end{definition}

\begin{definition}[Paired Boolean output]
\lean{Complexity.BinArith.pairO}
\label{Complexity.BinArith.pairO}
\leanok
Let $\sigma$ be a type and let $o : \sigma \times \{0,1\} \times \{0,1\} \to \{0,1\}$.
Define an output that is absent when no bit is pending, and otherwise returns the
Boolean value $o(s,x,y)$ from the state $s$, pending bit $x$, and input bit $y$.
\end{definition}

\begin{theorem}[Transduction of paired bits]
\lean{Complexity.BinArith.transduce_pairs}
\label{Complexity.BinArith.transduce_pairs}
\leanok
\uses{Complexity.BinArith.pairO, Complexity.BinArith.pairδ, Complexity.BinArith.serialOut, Complexity.transduce}
\difficulty{4}
Let $\sigma$ be a set, let $\delta : \sigma \times \{0,1\} \times \{0,1\} \to \sigma$
and $o : \sigma \times \{0,1\} \times \{0,1\} \to \{0,1\}$, let $s \in \sigma$, and let
$ps$ be a finite list of bit pairs. A one-pass transducer started in state $(\bot,s)$,
which stores the first bit of each pair and updates its state by $\delta$ on the second
bit while emitting $o$ of the current state and the pair, outputs exactly the serial
output obtained by processing the pairs in $ps$ from state $s$. Its input is the word
formed by listing each pair's two bits consecutively.
\end{theorem}

\begin{theorem}[Flattening a padded zip]
\lean{Complexity.BinArith.flatMap_zipChunk}
\label{Complexity.BinArith.flatMap_zipChunk}
\leanok
\uses{Complexity.BinArith.padZip}
\difficulty{5}
Let $a,b$ be finite bit strings and let $N$ be a natural number. For each
$i=0,\ldots,N-1$, form the pair consisting of the $i$th bit of each string, using $0$
when that position is beyond the end of the string. Listing the two entries of each pair
in order gives the same bit string as listing, for each $i=0,\ldots,N-1$, the $i$th bit
of $a$ followed by the $i$th bit of $b$, again using $0$ beyond either string's end.
\end{theorem}

\begin{definition}[Pair encoding after dropping leading bits]
\lean{Complexity.BinArith.zipStep}
\label{Complexity.BinArith.zipStep}
\leanok
\uses{Complexity.pairFstD, Complexity.pairSndD, Turing.pairEncode}
For a Boolean list $w$, take the two components obtained by decoding $w$ as a pair,
remove the first bit from each component, and pair-encode the resulting lists.
\end{definition}

\begin{theorem}[Iteration of the zip step]
\lean{Complexity.BinArith.iterate_zipStep}
\label{Complexity.BinArith.iterate_zipStep}
\leanok
\uses{Complexity.BinArith.zipStep, Complexity.pairFstD, Complexity.pairSndD, Turing.pairEncode}
\difficulty{4}
For every bit string $w$ and natural number $i$, applying the zip step $i+1$ times to
$w$ gives the pair encoding of the tails obtained by removing the first $i+1$ bits from
the total first and total second projections of $w$.
\end{theorem}

\begin{theorem}[Polynomial-time padded pair interleaving]
\lean{Complexity.BinArith.polyTimeComputable_zipStream}
\label{Complexity.BinArith.polyTimeComputable_zipStream}
\leanok
\uses{Complexity.BinArith.flatMap_zipChunk, Complexity.BinArith.iterate_zipStep, Complexity.BinArith.padZip, Complexity.BinArith.zipStep, Complexity.PolyTimeComputable, Complexity.PolyTimeComputable.append, Complexity.PolyTimeComputable.comp, Complexity.PolyTimeComputable.pairEncode, Complexity.length_pair_components_le, Complexity.pairFstD, Complexity.pairSndD, Complexity.polyTimeComputable_emitIter, Complexity.polyTimeComputable_headD, Complexity.polyTimeComputable_pairFstD, Complexity.polyTimeComputable_pairSndD, Complexity.polyTimeComputable_tail, Turing.length_pairEncode}
\difficulty{5}
For a bit string $z$, take the total first and second projections of $z$, paddedly pair
their bits for $|z|+2$ positions by using $0$ whenever either string has ended, and
concatenate each resulting pair in order. The function sending $z$ to this bit string is
polynomial-time computable.
\end{theorem}

\begin{theorem}[Finite-state serial output is polynomial-time computable]
\lean{Complexity.BinArith.polyTimeComputable_serial}
\label{Complexity.BinArith.polyTimeComputable_serial}
\leanok
\uses{Complexity.BinArith.padZip, Complexity.BinArith.pairO, Complexity.BinArith.pairδ, Complexity.BinArith.polyTimeComputable_zipStream, Complexity.BinArith.serialOut, Complexity.BinArith.transduce_pairs, Complexity.PolyTimeComputable, Complexity.PolyTimeComputable.comp, Complexity.pairFstD, Complexity.pairSndD, Complexity.polyTimeComputable_transduce}
\difficulty{3}
Let $\sigma$ be a finite set with decidable equality, let
$\delta:\sigma\times\{0,1\}^2\to\sigma$ be a state transition function, let
$o:\sigma\times\{0,1\}^2\to\{0,1\}$ be an output function, and fix an initial state
$s\in\sigma$. For each bit string $z$, decode its two components, pad them on the right
with zeros as needed, and form $|z|+2$ ordered pairs of bits. Starting from $s$, apply
$o$ to the current state and each successive bit pair to produce the output bit, then
update the state using $\delta$; the resulting bit-string function is polynomial-time
computable.
\end{theorem}
